\ysrc{27}
\section{粗模空间及其应用}

本节研究粗模空间 $Z^1(W_F^0/P,\hatG)\sslash\hatG$，即仿射概形范畴中的商，也称 GIT 商。具体地说，它是不变函数的谱：若 $Z^1(W_F^0/P,\hatG)=\Spec R_{\hatG,1_P}$，则
\[
Z^1(W_F^0/P,\hatG)\sslash\hatG=\Spec\bigl(R_{\hatG,1_P}\bigr)^{\hatG}.
\]

\subsection{几何不变量理论}
设 $\Lambda$ 仍如前述是系数环，$X=\Spec R$ 是带约化群 $\Lambda$-概形 $H$ 之作用的 $\Lambda$-概形。令 $X\sslash H:=\Spec R^H$。这里 $H$-不变量取群概形表示意义下的含义：作用对应一个余模结构 $R\xrightarrow{\rho}R\otimes_\Lambda\mathcal O(H)$，而 $R^H$ 是 $\mathrm{Id}\otimes1$ 与 $\rho$ 的等化子。显然 $X\sslash H$ 具有仿射 $\Lambda$-概形范畴中商所应有的泛性质。

关于这些对象有如下标准事实（回忆 $H$ 在此是约化群）。

\begin{enumerate}[leftmargin=2.4em,itemsep=0.15em]
\item $R^H$ 是有限生成 $\Lambda$-代数（$\Lambda$ 为特征 $0$ 的域时归功于 Mumford，正特征时归功于 Haboush，$\Lambda=\Z[\frac1p]$ 时归功于 Seshadri 与 Thomason [Tho87]）。
\item $X\sslash H$ 的构造与平坦基变换交换，但一般不与任意基变换交换。不过，若 $L$ 是 $\Lambda$ 上的域，则典范映射 $X_L\sslash H_L\to(X\sslash H)_L$ 是泛同胚 [Alp14]。
\item 设 $\Lambda=L$ 为代数闭域，$\pi:X\to X\sslash H$ 为典范映射。则对任意 $x\in(X\sslash H)(L)$，纤维 $\pi^{-1}(x)\subset X(L)$ 含有唯一的闭 $H(L)$-轨道。这给出 $X(L)$ 中的闭轨道与 $X\sslash H$ 的 $L$-点之间的双射。
\item Hilbert--Mumford 判据断言：点 $x\in X(L)$ 的轨道闭，当且仅当对任意余特征 $\lambda:\mathbb G_m\to H$，若态射 $t\mapsto\lambda(t)\cdot x$ 可延拓到 $\mathbb A^1$，则该延拓在 $t=0$ 处的值（通常记作 $\lim_{t\to0}\lambda(t)\cdot v$）属于 $x$ 的轨道。
\item 设 $X$ 既约，$Y$ 是 $X$ 的 $H$-稳定闭子概形，且对 $\Lambda$ 上任意代数闭域，$Y(L)$ 包含 $X(L)$ 的全部闭 $H(L)$-轨道。则自然映射 $Y\sslash H\to X\sslash H$ 是泛同胚，且在张量 $\Q$ 之后是同构。[这可由 [Alp14] 推出，说明见 [DHKM20, Prop. 4.17]]。
\end{enumerate}

\begin{exam}{5.1（Chevalley--Steinberg）}
令 $X=H$，$H$ 以共轭作用在自身上。设 $H$ 分裂，$T$ 是极大环面，正规化子为 $N$，Weyl 群为 $\Omega$。则 Chevalley--Steinberg 定理说：映射 $T\subset H$ 诱导出同构
\[
T\sslash N=T\sslash\Omega\xrightarrow{\ \sim\ }H\sslash H.
\]
它在 $L$-点上是双射，这一点容易由上面第 (4) 条中的 Hilbert--Mumford 判据推出。然而要证明它在任意环 $\Lambda$ 上成立，则需要更多关于不变量代数的工作，并涉及根系的组合学。此外，这也是一个取不变量恰好与任意基变换交换的例子。
\end{exam}

\begin{exam}{5.2（扭 Chevalley--Steinberg）}
令 $X=H$，$H$ 以 $\theta$-共轭作用在自身上：$(g,h)\mapsto gh\theta(g)^{-1}$，其中 $\theta$ 是某个自同构。等价地，我们考虑的是 $H$ 在陪集 $H\rtimes\theta$ 上的普通共轭作用。设 $\theta$ 正规化 $H$ 的 Borel 对 $(B,T)$（在代数闭域上总可化归为此情形）。特别地，$\theta$ 稳定 $N$ 并作用在 $\Omega=N/T$ 上。
\ysrc{28}
令 $N_\theta\subset H$ 是 $T\rtimes\theta$ 在 $H\rtimes\theta$ 中的稳定化子。它也是 $N$ 中不动点子群 $\Omega^\theta$ 的原像。扭 Chevalley--Steinberg 定理说：映射 $T\subset H$ 诱导出同构
\[
T\rtimes\theta\sslash N_\theta=T_\theta\sslash\Omega^\theta\xrightarrow{\ \sim\ }H\rtimes\theta\sslash H.
\]
这里 $T_\theta=T/(\mathrm{Id}-\theta)T$ 是 $T$ 的 $\theta$-余不变商环面。
\end{exam}

\subsection{抛物子群}
如前几节，$H$ 仍表示一个「广义约化」群概形，而所有作用最终都来自 $H^0$ 在 $H$ 上的共轭作用。给定定义在 $\Lambda$ 上的余特征 $\lambda:\mathbb G_m\to H^0$，存在光滑闭子群概形 $P_\lambda\subset H$，其点是那些 $x\in H$，使得轨道态射 $\mathbb G_m\to H$，$z\mapsto\lambda(z)x\lambda(z)^{-1}$ 可延拓到 $\mathbb A^1$。此外还有投影态射 $P_\lambda\xrightarrow{\pi_\lambda}M_\lambda:=C_H(\lambda)$，它把 $x$ 送到 $\lim_{t\to0}\lambda(t)x\lambda(t)^{-1}$。这样的子群概形 $P_\lambda$（相应地 $M_\lambda$）称为 $H$ 的 R-抛物子群（相应地 R-Levi 子群）。它们的意义来自 Hilbert--Mumford 判据。

\begin{rem}{5.3}
交 $P_\lambda\cap H^0$ 是 $H^0$ 的抛物子群概形，而 $M_\lambda\cap H^0$ 是 $P_\lambda\cap H^0$ 的 Levi 分支。此外，已知 $H^0$ 的任意抛物子群或 Levi 子群都具有这种形式。因此上述构造为不连通的约化群概形提供了一种抛物子群/Levi 子群的概念，例如 Richardson 在 [Ric88] 中用的就是这一概念。不过文献中还出现过别的概念。例如，对 $L$-群 $H=\hatG\rtimes W$ 的情形，Borel 在 [Bor79] 中把 $H$ 的抛物子群定义为 $H^0$ 的某个抛物子群的正规化子，且该抛物子群要满射到 $\pi_0(H)=W$ 上。在 [DM94] 中，Digne 与 Michel 把 $H$ 的抛物子群定义为 $H^0$ 的某个抛物子群的正规化子。在代数闭域上，Borel 意义下的抛物子群是 DM 意义下的抛物子群，而后者又是 Richardson 意义下（简称「R-」）的抛物子群，但两个反向蕴含都可能不成立。
\end{rem}

\begin{exam}{5.4}
$H$ 中 Borel 子群的正规化子是 R、B 或 DM 意义下的抛物子群。类似地，Borel 对的正规化子在每种意义下都是 Levi 子群。
\end{exam}

\subsection{半单同态}
设 $\Gamma$ 是抽象有限生成群，$H$ 是系数环 $\Lambda$ 上的广义约化群概形。我们关心 $H^0$ 通过共轭在 $X:=\Hom(\Gamma,H)$ 上的作用。

\begin{thm}{5.5}\label{thm:dat-5-5}
设 $L$ 是 $\Lambda$ 上的代数闭域，$\varphi:\Gamma\to H(L)$ 是同态。则下列性质等价：
\begin{enumerate}[leftmargin=2.2em,itemsep=0.1em,topsep=0.2em]
\item $\varphi$ 在 $H^0(L)$ 下的轨道在 $X(L)$ 中闭。
\item 对任意满足 $\varphi(\Gamma)\subset P(L)$ 的 R-抛物子群 $P$，都存在 $P$ 的 R-Levi 分支 $M$ 使 $\varphi(\Gamma)\subset M(L)$。
\item 对中心化子 $C_H(\varphi)$ 的任意极大环面 $S$，像 $\varphi(\Gamma)$ 不包含在 $C_H(S)$ 的真 R-抛物子群中。
\end{enumerate}
若这些性质成立，我们就说 $\varphi$ 是半单的。此时连通中心化子 $C_H(\varphi)^0$ 与连通正规化子 $N_H(\varphi(\Gamma))^0$ 都是约化群。
\end{thm}

\begin{proof}
性质 (2) 说的是 $\varphi(\Gamma)$ 在 Serre 意义下是 $H$-完全可约的。由 [BMR05, Thm. 3.1 与 \S6]，这等价于性质 (3)，而后者说的是 Zariski 闭包 $\overline{\varphi(\Gamma)}$ 是 $H$ 在 [Ric88] 意义下的强约化子群。
\ysrc{29}
现在取 $\Gamma$ 的生成元 $\gamma_1,\dots,\gamma_n$，于是 $X$ 是 $H^n$ 的闭子概形。由 [Ric88, Thm. 16.7]，性质 (3) 等价于：元素 $(\varphi(\gamma_1),\dots,\varphi(\gamma_n))\in H(L)^n$ 在 $H$ 的对角共轭作用下的 $H$-轨道 $\mathcal O$ 在 $H^n$ 中闭。但 $\varphi$ 的 $H$-轨道就是 $\mathcal O\cap X$，因而在 $X$ 中闭。最后一句陈述是 [BMR05] 的命题 3.12。
\end{proof}

\begin{exam}{5.6}
若 $\Gamma$ 有限且其阶在 $L$ 中可逆，则任意 $\varphi:\Gamma\to H(L)$ 都是半单的，因为它的 $H^0$-轨道既闭又开（如我们已看到的）。
\end{exam}

\begin{exam}{5.7}
设 $\Gamma=s^{\Z}$，$\sigma:=\varphi(s)$。若 $\sigma$ 是 $H(L)$ 的半单元，则 $\varphi$ 在上述意义下是半单的。若 $\pi_0(H)$ 的阶在 $L$ 中可逆，则逆命题成立；但一般而言可能不成立，例如 $H=\Z/\ell\Z$ 而 $L$ 的特征为 $\ell$。一般地，若 $\varphi$ 半单，则 $H^0$ 的每个被 $\sigma$ 正规化的 Borel 子群都含有被 $\sigma$ 正规化的极大环面（注意：由 Steinberg 的一个结果，$H(L)$ 中任意元素都正规化 $H^0$ 的某个 Borel 子群）。
\end{exam}

\begin{exam}{5.8}
若 $H=\mathrm{GL}_n$，则 $\varphi$ 作为同态半单，当且仅当对应的 $n$ 维表示在通常意义下半单。
\end{exam}

众所周知，半单表示限制到正规子群上仍然半单。下面的结果是它的推广。

\begin{thm}{5.9 ([BMR05, Thm. 3.10])}
设 $\Delta$ 是 $\Gamma$ 的正规子群，$\varphi:\Gamma\to H(L)$ 半单。则 $\varphi|_\Delta$ 半单。
\end{thm}

下面的推论在我们的情形中会用到。

\begin{cor}{5.10}\label{cor:dat-5-10}
设 $\Gamma=\Delta\rtimes\Sigma$，$\varphi:\Gamma\to H(L)$ 是同态，且 $\varphi(\Delta)$ 有限。令 $H_{\varphi|\Delta}:=N_H(\varphi(\Delta))$。则下列条件等价：
\begin{enumerate}[leftmargin=2.2em,itemsep=0.1em,topsep=0.2em]
\item $\varphi$ 半单；
\item $\varphi|_\Delta$ 半单，且 $(\varphi|_\Sigma)^{H_{\varphi|\Delta}}:\Sigma\to H_{\varphi|\Delta}(L)$ 半单。
\end{enumerate}
\end{cor}

要注意，当 $\varphi$ 半单时，$\varphi|_\Sigma$ 作为同态 $\Sigma\to H(L)$ 未必半单。例如，$S_3=\Z/3\Z\rtimes\{\pm1\}$ 的二维不可约 $\F_2$-表示限制到 $\{\pm1\}$ 上就不半单。

\begin{proof}
为简化记号，我们直接把 $H(L)$ 写成 $H$，对 $H_L$ 的其他闭代数子群也照此办理。又简记 $\varphi:=\phi|_\Delta$。假设 $\phi(\Delta)$ 有限，这蕴含 $(H_\varphi)^0=C_{H^0}(\varphi)^0$。令
\[
\begin{aligned}
\Hom(\Gamma,H)_\phi&:=\{\theta\in\Hom(\Gamma,H),\ \theta|_\Delta=\phi\},\\
\Hom(\Gamma,H)_{[\phi]}&:=\{\theta\in\Hom(\Gamma,H),\ \theta|_\Delta\ \text{与}\ \phi\ \text{在}\ H^0\ \text{下共轭}\}.
\end{aligned}
\]
记 $\rho:\Hom(\Gamma,H)\to\Hom(\Delta,H)$ 为限制映射，$H^0\cdot\phi$ 为 $\phi$ 的 $H^0$-轨道，则 $\Hom(\Gamma,H)_{[\phi]}=\rho^{-1}(H^0\cdot\phi)$。特别地，我们看到
\[
\phi\ \text{半单}\ \Longrightarrow\ \Hom(\Gamma,H)_{[\phi]}\ \text{在}\ \Hom(\Gamma,H)\ \text{中闭}.
\]
另一方面，$\Hom(\Gamma,H)_\phi=\Hom(\Gamma,H_\phi)_\phi$，而限制映射 $\Hom(\Gamma,H_\phi)_\phi\to\Hom(\Sigma,H_\phi)$ 是开闭浸入（如同引理 4.4 中那样）。
\ysrc{30}
于是我们看到：$(\varphi|_\Sigma)^{H_\phi}$ 半单，当且仅当 $\varphi^{H_\phi}$ 半单，等价地，轨道 $C_{H^0}(\phi)\cdot\varphi$ 在 $\Hom(\Gamma,H)_\phi$ 中闭。利用 $\Hom(\Gamma,H)_{[\phi]}=H^0\times_{C_{H^0}(\phi)}\Hom(\Gamma,H)_\phi$，我们推出
\[
(\varphi|_\Sigma)^{H_\phi}\ \text{半单}\iff H^0\cdot\varphi\ \text{在}\ \Hom(\Gamma,H)_{[\phi]}\ \text{中闭}.
\]
至此我们已证明 (2)$\Rightarrow$(1)。另一方向由上一个定理推出。
\end{proof}

现在转向我们所关心的情形。

\begin{prop}{5.11}\label{prop:dat-5-11}
设 $\Gamma=W_t^0=\langle \mathrm{Fr},s\rangle$，$\varphi:W_t^0\to H(L)$ 半单。如前令 $F=\varphi(\mathrm{Fr})$，$\sigma=\varphi(s)$。则：
\begin{enumerate}[leftmargin=2.2em,itemsep=0.1em,topsep=0.2em]
\item $\sigma$ 正规化 $H^0$ 的某个 Borel 对，而 $F$ 正规化 $C_H(\sigma)^0$ 的某个 Borel 对；
\item $\sigma$ 的阶有限，且该阶界不依赖于 $L$。
\end{enumerate}
\end{prop}

\begin{proof}
(1) 由上一个推论应用于 $\Delta=s^{\Z[\frac1p]}$ 与 $\Sigma=\mathrm{Fr}^{\Z}$，并应用例 5.7 得到。至于 (2)，注意：既然 $\sigma$ 正规化 $H^0$ 的某个 Borel 对，元素 $\sigma^{|\pi_0(H)|}$ 属于该 Borel 对的环面，从而是半单的。但在命题 3.3 的证明中，我们已经对与自身 $q$ 次幂共轭的半单元的阶算出了一个只依赖于 $q$ 与 $H$ 的界。
\end{proof}

注意，(2) 特别蕴含：任何半单的 $\varphi:W_t^0\to H(L)$ 都唯一连续地延拓到 $W_t$。更精确且更一般地，我们有：

\begin{cor}{5.12}
设 $P\subset P_F$ 是 $W_F$ 中 distinguished 的开子群。则存在开子群 $I\subset I_F$，使得 $W_F^0/P$ 满射到 $W_F/I$ 上，且对任意代数闭域 $L$，每个半单同态 $\varphi:W_F^0/P\to H(L)$ 都经过 $W_F/I$ 分解。因此，由限制诱导的态射
\[
\Hom(W_F/I,H)\sslash H^0\longrightarrow\Hom(W_F^0/P,H)\sslash H^0
\]
是泛同胚，且在将标量扩张到 $\Q$ 之后是同构。
\end{cor}

\begin{proof}
这由归约到驯顺情形的过程以及 5.1 中的 (5) 推出。
\end{proof}

\begin{cor}{5.13}\label{cor:dat-5-13}
若 $L$ 的特征为 $0$，则 $\varphi:W_F^0\to H(L)$ 半单，当且仅当 $\varphi(I_F^0)$ 有限且 $\varphi(\mathrm{Fr})$ 半单。
\end{cor}

特别地，$Z^1(W_F^0,\hatG)\sslash\hatG$ 的 $\C$-点正是通常局部 Langlands 参数的「无穷小特征」。

\subsection{在 $\Z[\frac1p]$ 上的连通性}
此处系数为 $\Lambda=\Z[\frac1p]$，我们将给出定理 3.13 的一个不同于 [DHKM20, \S4.6] 的证明梗概。更一般地，我们从 $\Z[\frac1p]$ 上的约化群 $\hatG$ 出发，它带有一个保持 Borel 对 $(\hat B,\hat T)$ 的 $W_t^0$ 的有限作用。

我们首先解释为何 $Z^1:=Z^1(W_t^0,\hatG)$ 连通，这里作一个简化假设：$W_t^0$ 在 $\hatG$ 上的作用非分歧。我们用的第一条原理是：映射
\[
Z^1(W_t^0,\hatG)\longrightarrow Z^1(W_t^0,\hatG)\sslash\hatG
\]
\ysrc{31}
在 $\pi_0$ 上诱导双射，因为 $\hatG$ 连通。于是只需证明 $Z^1(W_t^0,\hatG)\sslash\hatG$ 连通。

为此，我们引入 $Z^1(W_t^0,\hatG)$ 的两个闭子概形 $A$ 与 $B$：
\begin{align*}
A&:=\bigl((N_{\hat G}(\hat T)\rtimes\mathrm{Fr})\times\hat T\bigr)\cap Z^1(W_t^0,\hatG)\subset(\hatG\rtimes\mathrm{Fr})\times\hatG,\\
B&:=(\hatG\rtimes\mathrm{Fr})\times\{1\}\subset Z^1(W_t^0,\hatG).
\end{align*}
命题 5.11 (1) 蕴含：对 $\Lambda$ 上任意代数闭域 $L$，$Z^1(W_t^0,\hatG(L))$ 的每个闭 $\hatG(L)$-轨道都在 $A(L)$ 中有代表元。因此复合映射 $A\to Z^1\sslash\hatG$ 是满射，从而只需证明 $A$ 的所有连通分支都被映入 $Z^1\sslash\hatG$ 的同一个分支。为此注意：$A$ 是 $N_{\hat G}(\hat T)$ 上的闭子群概形，落在 $N_{\hat G}(\hat T)\times\hat T$ 中。具体地说，$A$ 是 $N_{\hat G}(\hat T)$ 上同源 $t\mapsto n\,\mathrm{Fr}(t)n^{-1}t^{-q}$ 的核。特别地，$A$ 是 $N_{\hat G}(\hat T)$ 上的有限平坦对角化群概形，其阶与 $p$ 互素。由于 $p$ 是 $N_{\hat G}(\hat T)$ 的每个分支上唯一可逆的素数，我们推出 $\pi_0(A)\xrightarrow{\sim}\pi_0(N_{\hat G}(\hat T))=\Omega$，且 $A$ 的所有分支都与 $A\cap B=(N_{\hat G}(\hat T)\rtimes\mathrm{Fr})\times\{1\}$ 相交。但 $B$ 连通，故 $B$ 被映入 $Z^1\sslash\hatG$ 的某个连通分支 $C$。由此 $A$ 的所有分支也都被映入 $C$，因而 $Z^1\sslash\hatG$ 连通。

现在回到一般情形：$s$ 在 $\hatG$ 上的作用非平凡（但仍正规化 $(\hat B,\hat T)$）。我们如下修改 $A$ 与 $B$ 的定义：
\begin{align*}
A&:=(N_s\rtimes\mathrm{Fr})\times(\hat T^{s,0}\rtimes s)\cap Z^1(W_t^0,\hatG),\\
B&:=(\hatG^s\rtimes\mathrm{Fr})\times\{1\rtimes s\}\subset Z^1(W_t^0,\hatG).
\end{align*}
这里 $\hat T^{s,0}$ 表示不动点子群概形 $\hat T^s$ 的极大环面，而
\[
N_s=\{n\in N_{\hat G}(\hat T),\ n.s^q(n)^{-1}\in\hat T^{s,0}\}.
\]
由命题 5.11 (1) 再次可知 $A$ 满射到 $Z^1\sslash\hatG$ 上（细节见 [DHKM24] 中的引理 2.9）。在此情形，[DHKM24] 的引理 2.8 表明：考察 $N_s$-环面 $N_s\times T^{s,0}$ 的自同源 $t\mapsto n\,\mathrm{Fr}(t)n^{-1}t^{-q}$，$A$ 就是截面 $N_s\to T^{s,0}$、$n\mapsto s^q(n)n^{-1}$ 在这个自同源下的原像。因此它仍然是某个有限平坦的、阶与 $p$ 互素的 $N_s$-群概形上的挠子。此外，loc.\ cit.\ 的引理 2.7 表明 $\pi_0(N_s)=\Omega^s$，且 $p$ 是 $N_s$ 各分支上唯一可逆的素数。于是如上，$\pi_0(A)\xrightarrow{\sim}\pi_0(N_s)=\Omega^s$。因此，要想按与上述简化情形相同的方式完成论证，只需知道
\begin{enumerate}[leftmargin=2.2em,itemsep=0.1em,topsep=0.2em]
\item 不动点子群概形 $\hatG^s$ 连通；
\item $\Omega^s$ 中的元素可以提升为 $A\cap B=N_{\hat G}(\hat T)^s$ 的一个点。
\end{enumerate}
若 $\hatG$ 单连通，则由 Steinberg 的一个著名结果，性质 (1) 成立。此外，[DHKM20] 的引理 4.27 与 4.28 通过 Langlands 的「z-扩张」技巧的某种变体，说明了如何归约到此情形。性质 (2) 一般并不成立，但若 $s$ 稳定某个固定数据（pinning），则由 [DM94, Thm. 1.15 iii)'] 它成立，因为此时 $s$ 是「拟中心的」。至此我们证明了定理 3.13。

\begin{rem}{5.14}
实际上，对上述论证而言，只需能把 $\Omega^s$ 中的元素提升为 $A\cap B$ 的某个 $\overline{\F}_\ell$-点（$\ell$ 为某个素数），而为此只需 $s$ 稳定 $\hatG_{\overline{\F}_\ell}$ 的某个固定数据。Cotner 的一个观察是：当 $\hatG$ 的中心在反转 $\ell$ 之后变光滑时，可以归约到此情形。
\end{rem}

\ysrc{32}
这足以证明他的定理 3.12 在 $\hatG$ 没有 A 型分支时的情形。另一方面，A 型单群的情形需要细致的研究和更复杂的论证。

\subsection{在 $\Z_\ell$ 上的连通性}
此处系数为 $\Z_\ell$（不久后为 $\overline{\F}_\ell$），我们将证明定理 4.12。我们沿用上一小节的记号，并进一步假设 $s$ 在 $\hatG$ 上的作用的阶是 $\ell$ 的幂。

由于 $Z^1(W_t^0,\hatG)_{\Z_\ell,1}$ 的所有分支都忠实平坦于 $\Z_\ell$ 上，只需证明特殊纤维 $Z^1_{\overline{\F}_\ell,1}=Z^1(W_t^0,\hatG)_{\overline{\F}_\ell,1}$ 连通，或等价地，商 $Z^1_{\overline{\F}_\ell,1}\sslash\hatG_{\overline{\F}_\ell}$ 连通。与上一小节一样，我们有满射 $A_{\overline{\F}_\ell}\to Z^1_{\overline{\F}_\ell}\sslash\hatG_{\overline{\F}_\ell}$，它诱导出满射 $A_{\overline{\F}_\ell,1}\to Z^1_{\overline{\F}_\ell,1}\sslash\hatG_{\overline{\F}_\ell}$，其中下标 $1$ 表示 $\sigma$ 的阶为 $\ell$-幂的那个轨迹。但 $\hat T^{s,0}$ 是环面，故 $\hat T^{s,0}(\overline{\F}_\ell)$ 没有 $\ell$-子群，从而 $A_{\overline{\F}_\ell,1}\subset B_{\overline{\F}_\ell}$。此外，由于 $s$ 的阶为 $\ell$-幂，$(\hatG_{\overline{\F}_\ell})^s$ 总是连通的，与它是否单连通无关 [DM94, Cor. 1.33]。于是 $B_{\overline{\F}_\ell}\simeq\hatG^s_{\overline{\F}_\ell}\rtimes\mathrm{Fr}$ 连通，从而 $Z^1_{\overline{\F}_\ell,1}\sslash\hatG_{\overline{\F}_\ell}$ 也连通。

其实我们还可以把这一论证再推进一步，得到 $Z^1_{\overline{\F}_\ell,1}\sslash\hatG_{\overline{\F}_\ell}$ 在同胚意义下的显式描述。具体地说，由于 $B_{\overline{\F}_\ell}$ 在 $\hatG^s_{\overline{\F}_\ell}$ 下稳定，包含 $B\subset Z^1$ 诱导出态射
\[
B_{\overline{\F}_\ell}\sslash\hatG^s_{\overline{\F}_\ell}
=(\hatG^s_{\overline{\F}_\ell}\rtimes\mathrm{Fr})\sslash\hatG^s_{\overline{\F}_\ell}
=Z^1(\langle\mathrm{Fr}\rangle,\hatG^s_{\overline{\F}_\ell})\sslash\hatG^s_{\overline{\F}_\ell}
\longrightarrow Z^1(W_t^0,\hatG)_{\overline{\F}_\ell,1}\sslash\hatG_{\overline{\F}_\ell}.
\]
我们已经知道这个态射是满射；它也是单射，因为把 $B(\overline{\F}_\ell)$ 的一个元素共轭到另一个元素的 $\hatG(\overline{\F}_\ell)$ 中元素必须属于 $\hatG(\overline{\F}_\ell)^s$。利用扭 Chevalley--Steinberg 定理，我们得到一个双射态射
\begin{equation}
(\hat T^s_{\overline{\F}_\ell})_{\mathrm{Fr}}\sslash\Omega^{s,\mathrm{Fr}}
=(\hat T^s_{\overline{\F}_\ell}\rtimes\mathrm{Fr})\sslash(N^s_{\overline{\F}_\ell})_{\mathrm{Fr}}
\longrightarrow Z^1(W_t^0,\hatG)_{\overline{\F}_\ell,1}\sslash\hatG_{\overline{\F}_\ell}.
\tag{5.1}
\end{equation}
这个态射也是有限的，这可由下面的定理 5.17 或定理 5.18 得到，故它是同胚（另一论证见 [DHKM20, Prop. 4.19 (2)] 的证明）。

\subsection{在 $\overline{\F}_\ell$ 上同胚意义下 $Z^1(W_F^0,\hatG)\sslash\hatG$ 的描述}
设 $G$ 是 $F$ 上的约化群。等式 (4.4) 给出分解
\[
(Z^1(W_F^0,\hatG)\sslash\hatG)_{\Z_\ell}
=\coprod_{\phi_\ell\in\Phi_\ell}\ \coprod_{\phi\in\Sigma_{\phi_\ell}}
Z^1_{\mathrm{Ad}_\phi}(W_t^0,C_{\hat G}(\phi_\ell)^0)_1\sslash C_{\hat G}(\phi_\ell)_\phi.
\]
这就把问题归约为：对保持 $\hatG$ 的 Borel 对、且驯顺 $\ell$-分歧的 $W_t$ 的有限作用，描述 $Z^1(W_t^0,\hatG)_{\overline{\F}_\ell,1}\sslash\hatG_{\overline{\F}_\ell}$——而这正是我们刚刚所做的。正如 [DHKM20, Cor. 4.21] 所说明的，对分支的一种可能更好的指标化方式是考虑集合 $\Psi(\overline{\F}_\ell)$，它是下列数对的 $\hatG(\overline{\F}_\ell)$-共轭类的代表元之集：
\begin{itemize}[leftmargin=1.8em,itemsep=0.1em]
\item 半单连续 $1$-上闭链 $\phi:I_F\to\hatG(\overline{\F}_\ell)$；
\item 元素 $\beta\in\mathrm{Trans}_{\hat G\rtimes\mathrm{Fr}}(\phi,\mathrm{Fr}\phi)/C_{\hat G}(\phi)^0$。
\end{itemize}
这里 $\mathrm{Fr}$ 表示 $W_F$ 中 Frobenius 的一个提升。注意：对 $\beta$ 在 $\mathrm{Trans}_{\hat G\rtimes\mathrm{Fr}}(\phi,\mathrm{Fr}\phi)$ 中的任意提升 $\tilde\beta$，存在唯一的连续 $1$-上闭链 $\varphi_{\tilde\beta}:W_F\to\hatG(\overline{\F}_\ell)$，它延拓 $\phi$ 且满足 ${}^L\varphi_{\tilde\beta}(\mathrm{Fr})=\tilde\beta$。

\ysrc{33}
对 $\Psi(\overline{\F}_\ell)$ 中的每个数对 $(\phi,\beta)$，取 $C_{\hat G}(\phi)^0$ 中的一个 Borel 对 $(\hat B_\phi,\hat T_\phi)$，以及一个正规化该 Borel 对的提升 $\tilde\beta$。于是我们得到映射 $\hat T_\phi\to Z^1(W_F^0,\hatG)_{\overline{\F}_\ell}$，$t\mapsto\varphi_{t\tilde\beta}$。利用 (5.1) 与分解 (4.4)，我们得到：

\begin{thm}{5.15}
上述所有映射的汇集诱导出同胚
\[
\coprod_{(\phi,\beta)\in\Psi(\overline{\F}_\ell)}(\hat T_\phi)_\beta\sslash(\Omega_\phi)^\beta
\xrightarrow{\ \sim\ }Z^1(W_F^0,\hatG)_{\overline{\F}_\ell}\sslash\hatG_{\overline{\F}_\ell}.
\]
\end{thm}

在此陈述中，$(\hat T_\phi)_\beta$ 表示 $\hat T_\phi$ 的 $\tilde\beta$-余不变量（它不依赖于 $\tilde\beta$ 的选取），而 $\Omega_\phi$ 表示 $\hat T_\phi$ 在（可能非连通的）完全中心化子 $C_{\hat G}(\phi)$ 中的 Weyl 群。

\begin{rem}{5.16}
就目前的状态而言，表示论仍然无法描述 $\Rep_{\overline{\F}_\ell}(G(F))$ 的 Bernstein 中心。然而，Fargues--Scholze 态射 [FS21, Thm. IX.5.2] 与上述拓扑描述一起，提示 $\Rep_{\overline{\F}_\ell}(G(F))$ 的既约 Bernstein 中心具有与 $\Rep_\C(G(F))$ 的经典 Bernstein 中心相当类似的描述。
\end{rem}

\subsection{$\Hom$ 概形的一个有限性性质}
设 $\Gamma$ 是有限生成离散群，$H$ 是诺特基 $S$ 上的「广义」约化群概形（见 3.1 的意义），其中含有一个闭的广义约化子群概形 $H'$。闭浸入 $\Hom(\Gamma,H')\hookrightarrow\Hom(\Gamma,H)$ 诱导出态射 $\Hom(\Gamma,H')\sslash H'\to\Hom(\Gamma,H)\sslash H$，而后者很少是闭浸入。不过，Cotner 最近证明了下述定理，解决了本文笔记早先版本中的一个猜想。

\begin{thm}{5.17 ([Cot23])}
映射 $\Hom(\Gamma,H')\sslash H'\to\Hom(\Gamma,H)\sslash H$ 是有限态射。
\end{thm}

例如，$\Gamma=\Z$ 的情形由（可能扭的）Steinberg--Chevalley 定理得到。此外，当 $S$ 是代数闭域时，这一陈述本质上已由 Vinberg [Vin96] 在特征 $0$ 情形、由 Martin [Mar03] 在正特征情形证明。

我们在本文中关心的情形是 $S=\Spec(\Z[\frac1p])$ 与 $\Gamma=W_F^0/P$，其中 $P$ 是 $P_F$ 中在 $W_F$ 里 distinguished 的开子群。这一情形最早在 [DHKM24] 中证明，依据的是下面这个或许有独立意义的结果，以及 $\Gamma=\Z$ 的情形（扭 Steinberg--Chevalley）。

\begin{thm}{5.18}
设 $\mathrm{Fr}$ 是 $W_F^0/P$ 中 Frobenius 的任意提升。求值映射 $\varphi\mapsto\varphi(\mathrm{Fr})$ 诱导出有限态射
\[
\Hom(W_F^0/P,H)\sslash H^0\longrightarrow H\sslash H^0.
\]
\end{thm}

\begin{proof}
我们只给出证明梗概，细节请参看 [DHKM24, Section 2]。利用分解 (4.3)，可归约为证明求值映射 $\varphi\mapsto{}^L\varphi(\mathrm{Fr})$ 诱导出有限态射
\[
Z^1(W_t^0,\hatG)\sslash\hatG\longrightarrow(\hatG\rtimes\mathrm{Fr})\sslash\hatG,
\]
其中 $\hatG$ 是带 $W_t^0$ 的有限作用的约化群概形，该作用保持一个 Borel 对。在此情形，回忆 5.4 小节的记号，特别是闭子概形
\ysrc{34}
$A=(N_s\rtimes\mathrm{Fr})\times(\hat T^{s,0}\rtimes s)\cap Z^1(W_t^0,\hatG)$。于是我们有交换方块
\begin{center}
\begin{tikzcd}
Z^1(W_t^0,\hatG)\sslash\hatG \arrow[r] \arrow[d,two heads] & \hatG\rtimes\mathrm{Fr}\sslash\hatG \arrow[d] \\
A\sslash\hat T^{s,0} \arrow[r] & N_s\rtimes\mathrm{Fr}\sslash\hat T^{s,0}
\end{tikzcd}
\end{center}
其中下方映射有限，因为 $A$ 是 $N_s$ 上某个有限平坦群概形下的挠子；右方竖直映射有限，由 Chevalley--Steinberg 定理（见 loc.\ cit.\ 的引理 2.2）。由此复合映射 $A\sslash\hat T^{s,0}\to\hatG\rtimes\mathrm{Fr}\sslash\hatG$ 有限。另一方面，左方竖直映射是满射，而 $Z^1(W_t^0,\hatG)\sslash\hatG$ 既约，故 $\mathcal O(Z^1(W_t^0,\hatG))^{\hatG}$ 单射进 $\mathcal O(A)^{\hat T^{s,0}}$。由诺特性，$\mathcal O(Z^1(W_t^0,\hatG))^{\hatG}$ 因而是 $\mathcal O(\hatG\rtimes\mathrm{Fr})^{\hatG}$ 上的有限生成模，此即所求。
\end{proof}

\subsection{游历代数与 Fargues--Scholze 比较}
本小节先设 $H$ 是仿射基 $S$ 上的广义约化群。当 $\Gamma$ 是有限表现离散群时，选一个表现就能把 $\Hom(\Gamma,H)$ 看成光滑仿射概形的极限（等化子）。还有一种更典范的做法：考虑由所有数对 $(F,\gamma)$ 构成的范畴 $\mathrm{FF}_\Gamma$，其中 $F$ 是有限生成自由群，$\gamma:F\to\Gamma$ 是同态。于是有表现
\[
\operatorname*{colim}_{(F,\gamma)\in\mathrm{FF}_\Gamma}\mathcal O(\Hom(F,H))\xrightarrow{\ \sim\ }\mathcal O(\Hom(\Gamma,H)).
\]
对每个整数 $n$ 取 $n$ 个字母上的自由群 $F_n$，我们可以把这个余极限写成心理上更好的形式：
\[
\operatorname*{colim}_{(n,\,F_n\xrightarrow{\gamma}\Gamma)}\mathcal O(H^n)\xrightarrow{\ \sim\ }\mathcal O(\Hom(\Gamma,H)).
\]
这一表现对于给概形 $\Hom(\Gamma,H)$ 赋予导出结构特别有用。不过，正如我们已经提到的，这样的导出结构在我们关心的情形 $\Gamma=W_F^0/P$ 中是平凡的。在此情形引入它的主要动机，其实来自 Fargues 与 Scholze 的工作 [FS21, Ch. IX] 中 à la V. Lafforgue 的游历算子（excursion operators）形式体系。

注意，在两边取 $H$-不变量就得到映射
\begin{equation}
\operatorname*{colim}_{(n,\,F_n\xrightarrow{\gamma}\Gamma)}\mathcal O(H^n)^H\longrightarrow\mathcal O(\Hom(\Gamma,H))^H.
\tag{5.2}
\end{equation}
左端的谱有时称为 $H$ 的伪表示空间。当 $S$ 含有一个特征 $0$ 的域时，映射 (5.2) 是同构，因为在 $H$-模上取 $H$-不变量是正合函子，而上述 $H$-代数的余极限是筛滤的，从而与到 $H$-模的遗忘函子交换。一般地，它在谱上诱导泛同胚，但未必是同构。

现在设 $H={}^LG$ 是 $S=\Z[\frac1p]$ 上约化群 $G$（$G$ 定义在 $F$ 上）的 Langlands 对偶，$\Gamma=W_F^0/P$。在 (5.2) 中再加上与投影 ${}^LG\to W_F^0$ 的通常相容性条件，所得的变体称为 $G$ 的游历代数，我们记作 $\mathrm{Exc}(W_F^0/P,\hatG)$。于是映射 (5.2) 写作
\begin{equation}
\mathrm{Exc}(W_F^0/P,\hatG)\longrightarrow\mathcal O(Z^1(W_F^0/P,\hatG))^{\hatG}.
\tag{5.3}
\end{equation}

\ysrc{35}
在此框架下，[FS21] 的主要结果之一是构造了一族相容的映射（随 $P$ 变化）
\begin{equation}
\mathrm{Exc}(W_F^0/P,\hatG)_{\Z_\ell}\longrightarrow\mathfrak Z_{\Z_\ell}(G(F)),
\tag{5.4}
\end{equation}
它们满足一系列非常好的性质。这里右端如前表示 Bernstein 中心。这就给出了一个明确的动机：确定映射 (5.3) 何时是同构。Fargues 与 Scholze 在 [FS21] 第 VIII.5 节中以引人注目的方式处理了这个问题。特别地，他们证明了：

\begin{thm}{5.19 (Thm. VIII.3.6 of [FS21])}
只要 $\ell$ 不整除 $\pi_1(\hatG)_{\mathrm{tors}}$ 的阶，映射 (5.3) 就是同构。
\end{thm}

\subsection{Fargues--Scholze 映射对光滑表示有限性的一个应用}
此处 $G$ 是 $F$ 上的约化群。在 [DHKM24] 中，Fargues--Scholze 映射 (5.4) 被用来证明下述纯表示论的结果。

\begin{thm}{5.20}
对 $G$ 的任意开紧子群 $H$，Hecke 代数 $\overline{\Z}_\ell[H\backslash G/H]$ 是其中心上的有限生成模，而该中心是有限生成交换 $\overline{\Z}_\ell$-代数。\tn{底本此处作 $\overline{\Z}_\ell$，疑为 $\Z_\ell$ 之误。}
\end{thm}

这个结果有几个推论，例如光滑 $\overline{\Z}_\ell$-表示上抛物诱导与限制之间的所谓「第二伴随性」。这些推论以及 Hecke 代数的诺特性，在一些附加假设下（例如 $G$ 是典型群，或 $p$ 足够大）早已为人所知，并且是用纯表示论（包括类型论）得到的，见 [Dat09]。然而除了 Helm [Hel16] 对 $\mathrm{GL}_n$ 的工作之外，关于中心则一无所知。Fargues--Scholze 映射在此背景下之所以是有力的工具，是因为它们提供了 Bernstein 中心中大量丰富的元素。确实，尽管人们对这些映射一般所知不多，但「简单」的事实——它们在环面上已知、与抛物诱导相容、与非分歧扭化相容——已足以推出 Bernstein 中心「足够大」。

上述定理的证明请参看 [DHKM24]，它相当短。我们在此只强调，一个关键的输入是：当 $M$ 是 $G$ 的 Levi 子群时，自然代数映射 $\mathrm{Exc}(W_F^0,\hatG)\to\mathrm{Exc}(W_F^0,\hat M)$ 的有限性。这一有限性可直接由定理 5.18 推出，特别地，既不需要定理 5.19，也不需要定理 5.17。

\section*{参考文献}
\begin{refs}
\item[{[Alp14]}] Jarod Alper, \emph{Adequate moduli spaces and geometrically reductive group schemes}, Algebr. Geom. \textbf{1} (2014), no. 4, 489--531, DOI 10.14231/AG-2014-022. MR3272912
\item[{[BMR05]}] Michael Bate, Benjamin Martin, and Gerhard R\"ohrle, \emph{A geometric approach to complete reducibility}, Invent. Math. \textbf{161} (2005), no. 1, 177--218, DOI 10.1007/s00222-004-0425-9. MR2178661
\item[{[Bor79]}] A. Borel, \emph{Automorphic $L$-functions}, Automorphic forms, representations and $L$-functions (Proc. Sympos. Pure Math., Oregon State Univ., Corvallis, Ore., 1977), Proc. Sympos. Pure Math., XXXIII, Amer. Math. Soc., Providence, RI, 1979, pp. 27--61. MR546608
\item[{[BZCHN20]}] David Ben-Zvi, Harrison Chen, David Helm, and David Nadler, \emph{Coherent Springer theory and the categorical Deligne--Langlands correspondence}, Invent. Math. \textbf{235} (2024), no. 2, 255--344, DOI 10.1007/s00222-023-01224-2. MR4689369
\ysrc{36}
\item[{[Chi18]}] Gianmarco Chinello, \emph{Blocks of the category of smooth $\ell$-modular representations of $\mathrm{GL}(n,F)$ and its inner forms: reduction to level $0$}, Algebra Number Theory \textbf{12} (2018), no. 7, 1675--1713, DOI 10.2140/ant.2018.12.1675. MR3871507
\item[{[Con14]}] Brian Conrad, \emph{Reductive group schemes} (English, with English and French summaries), Autour des sch\'emas en groupes. Vol. I, Panor. Synth\`eses, vol. 42/43, Soc. Math. France, Paris, 2014, pp. 93--444. MR3362641
\item[{[Cot23]}] Sean Cotner, \emph{Morphisms of character varieties}, Int. Math. Res. Not. IMRN \textbf{16} (2024), 11540--11548, DOI 10.1093/imrn/rnae124. MR4789091
\item[{[Cot24]}] Sean Cotner, \emph{Connected components of the moduli space of $L$-parameters}, arXiv:2404.16716, 2024.
\item[{[Dat09]}] Jean-Fran\c{c}ois Dat, \emph{Finitude pour les repr\'esentations lisses de groupes $p$-adiques} (French, with English and French summaries), J. Inst. Math. Jussieu \textbf{8} (2009), no. 2, 261--333, DOI 10.1017/S1474748008000054. MR2485794
\item[{[Dat17]}] Jean-Fran\c{c}ois Dat, \emph{A functoriality principle for blocks of $p$-adic linear groups}, Around Langlands correspondences, Contemp. Math., vol. 691, Amer. Math. Soc., Providence, RI, 2017, pp. 103--131, DOI 10.1090/conm/691/13895. MR3666052
\item[{[Dat18]}] Jean-Fran\c{c}ois Dat, \emph{Equivalences of tame blocks for $p$-adic linear groups}, Math. Ann. \textbf{371} (2018), no. 1-2, 565--613, DOI 10.1007/s00208-018-1648-1. MR3788858
\item[{[DHKM20]}] Jean-Fran\c{c}ois Dat, David Helm, Robert Kurinczuk, and Gilbert Moss, \emph{Moduli of Langlands parameters}, J. Eur. Math. Soc. (JEMS) \textbf{27} (2025), no. 5, 1827--1927, DOI 10.4171/jems/1599. MR4889237
\item[{[DHKM24]}] Jean-Fran\c{c}ois Dat, David Helm, Robert Kurinczuk, and Gilbert Moss, \emph{Finiteness for Hecke algebras of $p$-adic groups}, J. Amer. Math. Soc. \textbf{37} (2024), no. 3, 929--949, DOI 10.1090/jams/1034. MR4736530
\item[{[DL22]}] Jean-Fran\c{c}ois Dat and Thomas Lanard, \emph{Depth zero representations over $\Z[\frac1p]$}, arXiv:2202.03982, 2022; to appear in Algebra and Number Theory.
\item[{[DM94]}] Fran\c{c}ois Digne and Jean Michel, \emph{Groupes r\'eductifs non connexes} (French, with English and French summaries), Ann. Sci. \'Ecole Norm. Sup. (4) \textbf{27} (1994), no. 3, 345--406. MR1272294
\item[{[EH14]}] Matthew Emerton and David Helm, \emph{The local Langlands correspondence for $\mathrm{GL}_n$ in families} (English, with English and French summaries), Ann. Sci. \'Ec. Norm. Sup\'er. (4) \textbf{47} (2014), no. 4, 655--722, DOI 10.24033/asens.2224. MR3250061
\item[{[FG12]}] Jason Fulman and Robert Guralnick, \emph{Bounds on the number and sizes of conjugacy classes in finite Chevalley groups with applications to derangements}, Trans. Amer. Math. Soc. \textbf{364} (2012), no. 6, 3023--3070, DOI 10.1090/S0002-9947-2012-05427-4. MR2888238
\item[{[Fin21]}] Jessica Fintzen, \emph{Types for tame $p$-adic groups}, Ann. of Math. (2) \textbf{193} (2021), no. 1, 303--346, DOI 10.4007/annals.2021.193.1.4. MR4199732
\item[{[FS21]}] Laurent Fargues and Peter Scholze, \emph{Geometrization of the local Langlands correspondence}, arXiv:2102.13459, 2021; to appear in Ast\'erisque.
\item[{[Hel16]}] David Helm, \emph{The Bernstein center of the category of smooth $W(k)[\mathrm{GL}_n(F)]$-modules}, Forum Math. Sigma \textbf{4} (2016), Paper No. e11, 98, DOI 10.1017/fms.2016.10. MR3508741
\item[{[Hel23]}] Eugen Hellmann, \emph{On the derived category of the Iwahori--Hecke algebra}, Compos. Math. \textbf{159} (2023), no. 5, 1042--1110, DOI 10.1112/s0010437x23007145. MR4586564
\item[{[Hel20b]}] David Helm, \emph{Curtis homomorphisms and the integral Bernstein center for $\mathrm{GL}_n$}, Algebra Number Theory \textbf{14} (2020), no. 10, 2607--2645, DOI 10.2140/ant.2020.14.2607. MR4190413
\item[{[HM18]}] David Helm and Gilbert Moss, \emph{Converse theorems and the local Langlands correspondence in families}, Invent. Math. \textbf{214} (2018), no. 2, 999--1022, DOI 10.1007/s00222-018-0816-y. MR3867634
\item[{[Iwa55]}] Kenkichi Iwasawa, \emph{On Galois groups of local fields}, Trans. Amer. Math. Soc. \textbf{80} (1955), 448--469, DOI 10.2307/1992998. MR75239
\item[{[Jan04]}] Jens Carsten Jantzen, \emph{Nilpotent orbits in representation theory}, Lie theory, Progr. Math., vol. 228, Birkh\"auser Boston, Boston, MA, 2004, pp. 1--211. MR2042689
\item[{[Kal19]}] Tasho Kaletha, \emph{Regular supercuspidal representations}, J. Amer. Math. Soc. \textbf{32} (2019), no. 4, 1071--1170, DOI 10.1090/jams/925. MR4013740
\ysrc{37}
\item[{[Lan18]}] Thomas Lanard, \emph{Sur les $\ell$-blocs de niveau z\'ero des groupes $p$-adiques} (French, with English summary), Compos. Math. \textbf{154} (2018), no. 7, 1473--1507, DOI 10.1112/s0010437x18007133. MR3826462
\item[{[Lan21]}] Thomas Lanard, \emph{Sur les $\ell$-blocs de niveau z\'ero des groupes $p$-adiques II} (French, with English and French summaries), Ann. Sci. \'Ec. Norm. Sup\'er. (4) \textbf{54} (2021), no. 3, 683--750, DOI 10.24033/asens.2468. MR4311097
\item[{[Mar03]}] Benjamin M. S. Martin, \emph{Reductive subgroups of reductive groups in nonzero characteristic}, J. Algebra \textbf{262} (2003), no. 2, 265--286, DOI 10.1016/S0021-8693(03)00189-3. MR1971039
\item[{[MIY22]}] Alexander Bertoloni Meli, Naoki Imai, and Alex Youcis, \emph{The Jacobson--Morozov morphism for Langlands parameters in the relative setting}, Int. Math. Res. Not. IMRN \textbf{6} (2024), 5100--5165, DOI 10.1093/imrn/rnad217. MR4721051
\item[{[Ric88]}] R. W. Richardson, \emph{Conjugacy classes of $n$-tuples in Lie algebras and algebraic groups}, Duke Math. J. \textbf{57} (1988), no. 1, 1--35, DOI 10.1215/S0012-7094-88-05701-8. MR952224
\item[{[Tho87]}] R. W. Thomason, \emph{Equivariant resolution, linearization, and Hilbert's fourteenth problem over arbitrary base schemes}, Adv. in Math. \textbf{65} (1987), no. 1, 16--34, DOI 10.1016/0001-8708(87)90016-8. MR893468
\item[{[TvdK10]}] Antoine Touz\'e and Wilberd van der Kallen, \emph{Bifunctor cohomology and cohomological finite generation for reductive groups}, Duke Math. J. \textbf{151} (2010), no. 2, 251--278, DOI 10.1215/00127094-2009-065. MR2598378
\item[{[Vin96]}] E. B. Vinberg, \emph{On invariants of a set of matrices}, J. Lie Theory \textbf{6} (1996), no. 2, 249--269. MR1424635
\item[{[Yu01]}] Jiu-Kang Yu, \emph{Construction of tame supercuspidal representations}, J. Amer. Math. Soc. \textbf{14} (2001), no. 3, 579--622, DOI 10.1090/S0894-0347-01-00363-0. MR1824988
\item[{[Zhu25]}] Xinwen Zhu, \emph{Coherent sheaves on the stack of Langlands parameters}, The Langlands program: II. Geometrization of the Langlands correspondence, Proc. Sympos. Pure Math., vol. 112.2, Amer. Math. Soc., Providence, RI, 2025.
\end{refs}

\bigskip
\begin{center}
{\small Institut de Math\'ematiques de Jussieu, Sorbonne Universit\'e, Universit\'e Paris Cit\'e, CNRS.\\[0.2em]
4 place Jussieu, 75005 Paris, France\\[0.2em]
\itshape Email address: \texttt{jean-francois.dat@imj-prg.fr}}
\end{center}
